% TEMPLATE-MILC-v3.tex - plantilla maestra UNICA de las guias MILC v3.
%
% La usan las cuatro familias de guias, cada una con su driver:
%   · Grados 6-11          -> scripts/build-guias-g11.py
%   · Bebras               -> scripts/build-guias-bebras.py
%   · Semillero            -> scripts/build-guias-semillero.py
%   · Territorio interior  -> scripts/build-guias-territorio-interior.py
%
% Antes habia tres copias casi identicas (~400 lineas iguales, 6 distintas):
% cada mejora habia que aplicarla tres veces, y en la practica no se hacia
% (el motor de recursos vivia solo en dos de ellas). Lo que varia entre
% familias esta parametrizado en 6 placeholders que cada driver llena
% (se nombran aqui SIN su sintaxis real: el builder reemplaza por texto plano,
% tambien dentro de los comentarios, y un valor multilinea rompe el preambulo):
%
%   HEADER_IZQ        encabezado izquierdo de pagina
%   HEADER_DER        encabezado derecho de pagina
%   PORTADA_ETIQUETA  rotulo sobre el numero grande (GRADO/MOMENTO/MODULO)
%   PORTADA_SERIE     linea de serie de la portada
%   PORTADA_ID        identificador de la guia en la portada
%   PORTADA_CIRCULO   el circulito de grado (°), o vacio si no aplica
%   PDF_TITULO        titulo en texto plano (sin \textbf ni comillas TeX) para
%                     los metadatos del PDF (/Title); lo produce lib_xelatex.texto_pdf
%
% Composicion (auditoria 2026-09-03): las bandas de estacion piden espacio
% (\Needspace*) y prohiben el salto justo despues (\nopagebreak), asi no quedan
% huerfanas al pie; las cajas largas (softbox, drawbox, infoband, puentebox) son
% `breakable`, asi una caja de media pagina ya no arrastra todo a la siguiente
% dejando la anterior al 40 %. Todo texto sobre mostaza o verde va en milcNegro
% (blanco sobre #E5B400 da 1,93:1 y sobre #58A923 2,95:1; ninguno pasa WCAG AA).
% El resto de claves las documenta content/guias/_SCHEMA.md. El script valida
% que no queden placeholders sin reemplazar antes de escribir el archivo final.

% Etiquetado PDF (latex-lab): /Lang, estructura y texto alternativo de las
% figuras, para lectores de pantalla. Debe ir ANTES de \documentclass.
\DocumentMetadata{lang=es-CO,tagging=on}
\documentclass[11pt,letterpaper]{article}
% headheight=24pt: la cabecera derecha lleva dos lineas (identidad de la guia
% y estacion actual). headsep se acorta en la misma medida para que el bloque
% de texto quede exactamente donde estaba con headheight=12pt.
\usepackage[margin=2.0cm,headheight=24pt,headsep=13pt]{geometry}
\usepackage[table]{xcolor}
\usepackage{fontspec}
\usepackage[spanish]{babel}
\usepackage{tikz}
% Librerías TikZ para los diagramas de las guías (bloque `recursos:`):
%  · babel  → babel en español vuelve activos caracteres como «>», y TikZ los usa
%             en sus opciones (>=stealth, ->). Sin esta librería, cualquier
%             diagrama con flechas revienta con «\language@active@arg> has an
%             extra }». Debe ir ANTES de que se dibuje nada.
%  · positioning        → sintaxis «below left=2pt and 3pt».
%  · decorations.pathreplacing → llaves ({brace}) para señalar rangos.
\usetikzlibrary{babel,positioning,decorations.pathreplacing}
\usepackage{graphicx}
% Circuitos: el trabajo de electrónica se hace hoy con micro:bit y sus sensores
% integrados (MakeCode, sin cableado), así que no se necesitan esquemáticos.
% Si algún día entran montajes con protoboard, basta descomentar esta línea:
% los diagramas .tex/.tikz declarados en `recursos:` ya se incrustan solos.
% \usepackage{circuitikz}
\usepackage[most]{tcolorbox}
\usepackage{tabularx}
\usepackage{array}
\usepackage{fancyhdr}
\usepackage{titlesec}
\usepackage[document]{ragged2e}  % bandera a la derecha en todo el cuerpo
\usepackage{enumitem}
\usepackage{microtype}
\usepackage{etoolbox}
\usepackage{needspace}
% hyperref al final del preambulo: metadatos del PDF (titulo, autor, idioma,
% familia), marcadores por estacion (\pdfbookmark) y enlaces sin recuadro.
\usepackage[hidelinks,
  pdftitle={Mi identidad digital — cómo me presento en la red sin perderme},
  pdfauthor={PhD. Álvaro Cárdenas Orozco},
  pdfsubject={Tecnología e Informática · Grado 6 · Guía 2 · MILC v3},
  pdflang={es-CO},
  bookmarksopen=true,bookmarksnumbered=false]{hyperref}

% Helvetica Neue no trae la flecha U+2192, asi que cada «→» escrito literalmente
% en el YAML se compilaba a NADA, en silencio (462 flechas en 61 guias: rutas del
% tipo «paleta → tipografia → lockup» salian rotas). Mapeamos el caracter a su
% version matematica, que si tiene glifo. Asi el autor puede seguir escribiendo
% «→» comodamente en el YAML.
\usepackage{newunicodechar}
\newunicodechar{→}{\ensuremath{\rightarrow}}
% Mismo problema con los simbolos de marca y vinetas: el «✓» de «una palomita ✓»
% salia como cuadrito vacio en el PDF de todas las guias que lo usaban.
\usepackage{pifont}
\newunicodechar{✓}{\ding{51}}
\newunicodechar{✗}{\ding{55}}
\newunicodechar{✔}{\ding{52}}
\newunicodechar{•}{\textbullet}
\newunicodechar{≥}{\ensuremath{\geq}}
\newunicodechar{≤}{\ensuremath{\leq}}

\setmainfont{Helvetica Neue}
\setsansfont{Helvetica Neue}
\setlength{\parindent}{0pt}
\setlength{\parskip}{6pt}
% Sin particion de palabras: en 6.º-7.º una palabra cortada por guion al final
% de linea («Comuni-cacion») frena la lectura mucho mas que una linea corta.
\hyphenpenalty=10000
\exhyphenpenalty=10000
\pretolerance=2000
\tolerance=3000
\setlength{\emergencystretch}{3em}
\linespread{1.10}
\renewcommand{\arraystretch}{1.25}
\setlist{leftmargin=5mm,itemsep=1mm,topsep=1mm}
\newcolumntype{Y}{>{\RaggedRight\arraybackslash}X}

\definecolor{milcMagenta}{HTML}{E6007E}
\definecolor{milcVino}{HTML}{5A0038}
\definecolor{milcTurquesa}{HTML}{0093A5}
\definecolor{milcVerde}{HTML}{58A923}
\definecolor{milcMostaza}{HTML}{E5B400}
\definecolor{milcOcre}{HTML}{B86600}
\definecolor{milcNegro}{HTML}{241816}
\definecolor{milcGris}{HTML}{F7F7F4}
\definecolor{milcLinea}{HTML}{E8E2DD}
\definecolor{milcGradePrimary}{HTML}{0066FF}
\definecolor{milcGradeDark}{HTML}{003D99}
\definecolor{milcGradeSoft}{HTML}{D6E8FF}
\definecolor{milcGradeAccent}{HTML}{CFE7FF}
\definecolor{milcGradeText}{HTML}{FFFFFF}
\color{milcNegro}

\pagestyle{fancy}
\fancyhf{}
% Cabecera de dos lineas: identidad de la guia arriba y, a la derecha y en
% gris, la estacion en curso (\markright desde \milcestacion). Antes de la
% primera estacion la segunda linea va vacia; el \strut mantiene la altura
% para que la linea de arriba no baile entre paginas.
\lhead{\small Tecnología e Informática · Grado 6\\[-1pt]{\footnotesize\strut}}
\rhead{\small Guía 2 · MILC v3\\[-1pt]{\footnotesize\strut\color{milcNegro!65}\rightmark}}
\cfoot{\small\thepage}
\renewcommand{\headrulewidth}{0pt}

% Sobre mostaza (#E5B400) y verde (#58A923) el texto va en milcNegro; sobre el
% resto de la paleta, en blanco. Se usa dentro de fonttitle/fontupper, donde un
% \color posterior manda sobre coltitle.
\newcommand{\milctintasobre}[1]{%
  \ifboolexpr{test{\ifstrequal{#1}{milcMostaza}} or test{\ifstrequal{#1}{milcVerde}}}%
    {\color{milcNegro}}{\color{white}}}

\newtcolorbox{titlebox}[1]{
  enhanced,colback=#1,colframe=#1,arc=7mm,boxrule=0pt,
  left=7mm,right=7mm,top=3mm,bottom=3mm,
  before skip=8pt,after skip=8pt,
  fontupper=\sffamily\bfseries\Large\color{white}
}

\newtcolorbox{softbox}[3]{
  enhanced,breakable,pad at break=3mm,
  colback=#2,colframe=#1,borderline west={4pt}{0pt}{#1},
  arc=2mm,boxrule=.4pt,left=4mm,right=4mm,top=3mm,bottom=3mm,
  before skip=6pt,after skip=8pt,title={#3},coltitle=white,
  colbacktitle=#1,fonttitle=\sffamily\bfseries\milctintasobre{#1},fontupper=\RaggedRight,
  attach boxed title to top left={xshift=3mm,yshift=-2mm},
  boxed title style={arc=2mm, boxrule=0pt}
}

\newtcolorbox{drawbox}[2]{
  enhanced,breakable,pad at break=4mm,
  colback=white,colframe=#1,arc=4mm,boxrule=1pt,
  left=5mm,right=5mm,top=4mm,bottom=4mm,before skip=8pt,after skip=8pt,
  title={#2},coltitle=white,colbacktitle=#1,fonttitle=\sffamily\bfseries\milctintasobre{#1},
  fontupper=\RaggedRight,
  attach boxed title to top left={xshift=4mm,yshift=-2mm},
  boxed title style={arc=3mm, boxrule=0pt}
}

\newtcolorbox{infoband}[1]{
  enhanced,breakable,pad at break=2mm,colback=#1!8,colframe=#1!50,arc=2mm,boxrule=.5pt,
  left=4mm,right=4mm,top=2mm,bottom=2mm,before skip=4pt,after skip=4pt,
  fontupper=\small\RaggedRight
}

% Puente narrativo entre secciones (contrato editorial Conéctate).
% Da continuidad: "Acabas de X. Ahora vas a Y porque Z."
% Visualmente: banda izquierda en color del grado, texto en itálica pequeña.
\newtcolorbox{puentebox}{
  enhanced,breakable,pad at break=2mm,colback=milcGradeSoft,colframe=milcGradeDark,
  borderline west={3pt}{0pt}{milcGradeDark},
  arc=0mm,boxrule=0pt,left=5mm,right=4mm,top=2.5mm,bottom=2.5mm,
  before skip=4pt,after skip=8pt,
  fontupper=\itshape\small\color{milcNegro}\RaggedRight
}
% La flecha va en modo matematico a proposito: Helvetica Neue NO tiene el glifo
% U+2192, asi que \textrightarrow se compilaba a NADA («Missing character: There
% is no → in font Helvetica Neue») y el puente salia sin flecha. En modo
% matematico la flecha viene de la fuente matematica y si se dibuja.
\newcommand{\puente}[1]{%
  \begin{puentebox}%
    \textcolor{milcGradeDark}{$\rightarrow$}\hspace{2mm}#1%
  \end{puentebox}%
}

\newcommand{\linead}[1]{\noindent\color{milcLinea}\rule{#1}{0.5pt}\color{milcNegro}}
\newcommand{\respuesta}[1]{\par\vspace{2mm}\foreach \n in {1,...,#1}{\noindent\color{milcLinea}\rule{\linewidth}{0.5pt}\par\vspace{5mm}}\color{milcNegro}}
\newcommand{\checkbox}{\textcolor{milcGradeDark}{\fbox{\phantom{x}}}\hspace{5pt}}

% ─── Listas aireadas para las guías socioemocionales ────────────────────
% Componer las verificaciones como «\checkbox texto\\» deja el texto que salta
% de línea debajo del cuadrito y apelmaza el bloque. Estos entornos alinean la
% continuación y airean los ítems. Los usa Territorio Interior; cualquier
% familia puede hacerlo.
\newlist{checklista}{itemize}{1}
\setlist[checklista]{
  label=\textcolor{milcGradeDark}{\fbox{\phantom{x}}},
  leftmargin=9mm, labelsep=3.2mm, itemsep=2.4mm,
  topsep=2.5mm, parsep=0pt, partopsep=0pt, before=\RaggedRight
}
\newlist{pasoslista}{enumerate}{1}
\setlist[pasoslista]{
  label=\textcolor{milcGradeDark}{\sffamily\bfseries\arabic*.},
  leftmargin=8mm, labelsep=3mm, itemsep=2.8mm,
  topsep=1mm, parsep=0pt, partopsep=0pt, before=\RaggedRight
}

\newcommand{\phasecard}[4]{
\begin{tcolorbox}[enhanced,colback=#3,colframe=#2,arc=3mm,boxrule=.5pt,height=3.05cm,valign=center,halign=center,left=1.5mm,right=1.5mm,top=2mm,bottom=2mm]
{\sffamily\bfseries\fontsize{22}{24}\selectfont\textcolor{#2}{#1}}\par
{\sffamily\bfseries\small #4}
\end{tcolorbox}
}

% ── Piezas del rediseno 2026 (legibilidad 6.º-11.º) ───────────────────
% Banda de estacion: reemplaza el titlebox macizo. Titulo a la izquierda,
% dato practico (tiempo/modalidad) a la derecha, en una sola linea.
% Nunca huerfana: pide 9 lineas libres (o salta de pagina rellenando con
% \vfill) y prohibe el salto justo despues de la banda. Nueve y no seis:
% la caja partible que sigue necesita ~80pt (titulo adosado, relleno y dos
% lineas) para arrancar en la misma pagina; con menos, tcolorbox la manda
% entera a la siguiente y la banda se queda sola. El penalty 10000 va
% en `after`, ANTES del espacio posterior: un `after skip` (aunque sea 0pt)
% es un \vskip tras la caja, y ese glue es un punto de corte legal que un
% \nopagebreak escrito despues de \end{tcolorbox} ya no cierra. Cada banda
% deja marcador PDF y actualiza la cabecera derecha.
\newcounter{milcestacion}
\newcommand{\milcestacion}[2]{%
\Needspace*{9\baselineskip}%
\stepcounter{milcestacion}%
\pdfbookmark[1]{#2}{milcestacion\themilcestacion}%
\markright{#2}%
\begin{tcolorbox}[enhanced,colback=#1,colframe=#1,arc=2mm,boxrule=0pt,
  left=5mm,right=5mm,top=2.6mm,bottom=2.6mm,before skip=11pt,
  after={\par\nopagebreak\vskip7pt}]
{\sffamily\bfseries\fontsize{15}{18}\selectfont\milctintasobre{#1}#2}
\end{tcolorbox}}

% Tarjeta de paso numerada: sustituye las tablas «Pilar / Aplicacion», que
% eran formato de consulta docente, no de estudiante. El numero va en un
% circulo sobre el borde para que la secuencia se lea de un vistazo.
\newcommand{\milcpaso}[3]{%
\Needspace{4\baselineskip}%
\begin{tcolorbox}[enhanced,colback=white,colframe=milcLinea,arc=1.5mm,boxrule=.9pt,
  left=14mm,right=4mm,top=3mm,bottom=3mm,before skip=4pt,after skip=4pt,
  fontupper=\RaggedRight,
  overlay={\node[anchor=north west,circle,fill=#2,text=white,inner sep=0pt,
    minimum size=7.4mm,font=\sffamily\bfseries\small]
    at ([xshift=3.6mm,yshift=-3.2mm]frame.north west) {#1};}]
#3
\end{tcolorbox}}

% Casilla marcable de tamano util con lapiz (la anterior era \fbox{\phantom{x}},
% de alto variable segun el texto que la seguia).
\newcommand{\milcmarca}{\framebox[3.8mm]{\rule{0pt}{3.4mm}}}

% Fila marcable de lista suelta (checks de la estacion 1).
\newcommand{\milccheckrow}[1]{%
\par\vspace{1.8mm}\noindent
\makebox[6.4mm][l]{\raisebox{-.55mm}{\textcolor{milcNegro}{\milcmarca}}}%
\parbox[t]{\dimexpr\linewidth-6.4mm\relax}{\RaggedRight #1}\par}

% Celda marcable dentro de la rubrica (tabla de autoevaluacion). Misma
% sangria francesa, pero pensada para vivir dentro de una columna X.
\newcommand{\milcopcion}[1]{%
\makebox[5.2mm][l]{\raisebox{-.55mm}{\textcolor{milcNegro}{\milcmarca}}}%
\parbox[t]{\dimexpr\linewidth-5.2mm\relax}{\RaggedRight #1}}

% ── Recursos visuales de la guía ──────────────────────────────────────
% Declarados en el bloque `recursos:` del YAML y emitidos por el builder.
% \guiaFigura   → imagen rasterizada o PDF (foto, carta celeste, montaje).
% \guiaDiagrama → fuente TikZ/circuitikz (.tex) incrustada como vector:
%                 es la vía para circuitos y esquemáticos editables.
% [#1] = texto alternativo (alt) · #2 = ruta relativa al .tex · #3 = pie
% El alt llega al PDF etiquetado (/Alt) y lo lee el lector de pantalla.
\newcommand{\guiaFigura}[3][]{%
\begin{center}
\includegraphics[width=0.88\linewidth,height=0.38\textheight,keepaspectratio,alt={#1}]{#2}\\[1.5mm]
{\footnotesize\itshape\textcolor{milcNegro!75}{#3}}
\end{center}
}

\newcommand{\guiaDiagrama}[2]{%
\begin{center}
\input{#1}\\[1.5mm]
{\footnotesize\itshape\textcolor{milcNegro!75}{#2}}
\end{center}
}

\begin{document}
\thispagestyle{empty}

% ── Portada ───────────────────────────────────────────────────────────
% Antes: un rectangulo de color a pagina completa. Se veia bien en pantalla,
% pero en la fotocopiadora del colegio se comia el toner y salia gris sucio.
% Ahora la banda ocupa el tercio superior y el resto va en blanco.
\begin{tikzpicture}[remember picture,overlay]
  \path[fill=milcGradePrimary]
    (current page.north west) rectangle ([yshift=-10.6cm]current page.north east);

  \node[anchor=north west,text=milcGradeText] at ([xshift=2.0cm,yshift=-1.55cm]current page.north west)
    {\sffamily\bfseries\fontsize{12}{14}\selectfont GRADO};
  \node[anchor=north west,text=milcGradeText,opacity=.85] at ([xshift=2.0cm,yshift=-2.35cm]current page.north west)
    {\sffamily\bfseries\fontsize{8.5}{10}\selectfont SERIE GUÍAS · TECNOLOGÍA E INFORMÁTICA};
  \node[anchor=north east,text=milcGradeText,opacity=.85,align=right] at ([xshift=-2.0cm,yshift=-1.55cm]current page.north east)
    {\sffamily\bfseries\fontsize{9}{12}\selectfont I.E. Sor María Juliana\\Cartago, Valle del Cauca};

  \node[anchor=west,text=milcGradeText] (gradonum) at ([xshift=1.85cm,yshift=-7.1cm]current page.north west)
    {\sffamily\bfseries\fontsize{112}{112}\selectfont 6};
  % El circulito de grado (°) solo aplica a las guias de grado («11°»); en
  % Bebras y Semillero el numero grande es un momento/modulo, no un grado.
  % Se posiciona relativo al nodo `gradonum`, no en coordenadas absolutas.
  \draw[milcGradeText,line width=1.6pt]
    ([xshift=2.0mm,yshift=-4.5mm]gradonum.north east) circle (.15cm);

  \node[anchor=west,text=milcGradeText,text width=11.4cm,align=left] at ([xshift=1.5cm,yshift=0.5cm]gradonum.east)
    {
     {\sffamily\bfseries\fontsize{9}{11}\selectfont\MakeUppercase{Período 1 · Comunicación e identidad digital}}\\[3mm]
     {\sffamily\bfseries\fontsize{21}{25}\selectfont\setlength{\baselineskip}{27pt}Mi identidad\\[4pt]digital ---\\[4pt]quién soy en la red}\\[3.5mm]
     {\sffamily\bfseries\fontsize{15}{18}\selectfont Guía 2-6-TIC}
    };
\end{tikzpicture}

\vspace*{9.4cm}
\vfill

{\sffamily\bfseries\fontsize{8}{10}\selectfont\textcolor{milcGradeDark}{LO QUE TE LLEVAS HOY}}

\begin{tcolorbox}[enhanced,colback=white,colframe=milcNegro,arc=2mm,boxrule=1.6pt,
  left=5mm,right=5mm,top=4mm,bottom=4mm,before skip=3pt,after skip=12pt,
  fontupper=\RaggedRight\fontsize{11}{15.5}\selectfont]
Una \textbf{tarjeta de identidad digital} hecha en el cuaderno, con tu nombre, tu foto-dibujo, 3 cosas que te gusta hacer, 2 cosas que cuidas (que NO le cuentas a desconocidos), y \textbf{1 frase que te describa que tú mismo escribas} (no la copies de internet). Producto físico, no digital todavía. Lo muestras al final de la clase.
\end{tcolorbox}

\begin{tcolorbox}[enhanced,colback=milcGris,colframe=milcGris,arc=2mm,boxrule=0pt,
  left=5mm,right=5mm,top=3.5mm,bottom=3.5mm,after skip=12pt]
{\sffamily\bfseries\fontsize{8}{10}\selectfont\textcolor{milcNegro!55}{ESTA GUÍA ES DE}}\\[3mm]
\begin{tabularx}{\linewidth}{X p{3.2cm} p{3.2cm}}
\linead{\linewidth} & \linead{3.0cm} & \linead{3.0cm}\\[-1mm]
{\fontsize{8.5}{10}\selectfont\textcolor{milcNegro!55}{Nombre completo}} &
{\fontsize{8.5}{10}\selectfont\textcolor{milcNegro!55}{Grupo}} &
{\fontsize{8.5}{10}\selectfont\textcolor{milcNegro!55}{Fecha}}\\
\end{tabularx}
\end{tcolorbox}

\vfill

\noindent\textcolor{milcLinea}{\rule{\linewidth}{1pt}}\\[2mm]
{\sffamily\bfseries\fontsize{9.5}{12}\selectfont Autor: Dr. Álvaro Cárdenas Orozco}\\[1mm]
{\sffamily\fontsize{8.5}{11}\selectfont\textcolor{milcNegro!55}{Metodología MILC · apertura ancestral · pensamiento computacional · triángulo Dussel–estoicismo–Floridi}}

\newpage

\section*{Apertura: Mi identidad digital --- cómo me presento en la red sin perderme}

\begin{softbox}{milcGradeDark}{milcGradeSoft}{Saber ancestral}
\textbf{¿Sabes qué hacía el pregonero del barrio?} Antes de los celulares, en cada cuadra del centro de Cartago caminaba el pregonero. Tocaba una campanita y gritaba: \emph{``Reunión del barrio el sábado, traer ladrillos para arreglar el parque.''} Llevaba mensajes casa a casa. Su voz lo identificaba: \textbf{ese pregonero era ese pregonero}. Si dañaba un mensaje, lo conocían y le decían: \emph{``Don Aurelio, esa noticia no la entregó bien.''} Si la entregaba bien, lo buscaban para más. \textbf{Su nombre y su voz eran su identidad}. La cuidaba porque la necesitaba para vivir. En internet pasa parecido: lo que firmas con tu nombre se queda. Si lo cuidas, te abre puertas. Si no, también se queda. Solo que en internet, mucha más gente te ve.
\end{softbox}

\begin{softbox}{milcTurquesa}{milcGris}{Saber contemporáneo}
Tu \textbf{identidad digital} es el grupo de cosas que la gente puede ver de ti cuando te buscan en la red. No es solo tu nombre. Es: \textbf{tu foto, tu apodo, lo que escribes, lo que comentas, los videos que subes, los grupos donde estás}. Cada vez que pones algo en internet, agregas un pedacito a esa identidad. Algunas cosas se borran fácil. Otras no. Por eso esta sesión es importante: \textbf{antes de poner algo en la red, decide tú qué quieres que esa identidad cuente de ti}. Empezamos hoy con un ejercicio muy concreto: vas a hacer una \emph{tarjeta de identidad digital} en tu cuaderno. Sin computador, sin celular. Solo lápiz.
\end{softbox}

\begin{softbox}{milcMostaza}{milcGris}{Pregunta-puente}
\textit{Si alguien te buscara en internet hoy mismo, ¿qué \textbf{cosas tuyas} encontraría? ¿Te gusta lo que vería?}
\end{softbox}

\begin{softbox}{milcVerde}{milcGris}{Saber y saber hacer}
Al terminar la clase: (1) podrás \textbf{decir con tus palabras} qué es identidad digital; (2) sabrás \textbf{distinguir} cosas que sí se cuentan en la red de cosas que no; (3) habrás \textbf{hecho} tu tarjeta de identidad digital en el cuaderno.
\end{softbox}



\small
\begin{tabularx}{\textwidth}{>{\bfseries}p{3.0cm}Y}
\rowcolor{milcGradePrimary!12}Autor & Dr. Álvaro Cárdenas Orozco\\
DBA & Reconoce principios y conceptos propios de la tecnología, relacionando artefactos con su utilización segura (MEN, T\&I 6°)\\
Referentes & Saber ancestral del pregonero · Pensamiento computacional inicial · Cuidado de la palabra\\
Producto final & Una \textbf{tarjeta de identidad digital} hecha en el cuaderno, con tu nombre, tu foto-dibujo, 3 cosas que te gusta hacer, 2 cosas que cuidas (que NO le cuentas a desconocidos), y \textbf{1 frase que te describa que tú mismo escribas} (no la copies de internet). Producto físico, no digital todavía. Lo muestras al final de la clase.\\
\end{tabularx}
\normalsize

\puente{Hoy vas a hacer 3 cosas: leer un poco, pensar un poco y dibujar tu tarjeta de identidad digital.}

\milcestacion{milcGradeDark}{Ruta MILC: así vamos a aprender}
\begin{tabularx}{\textwidth}{>{\centering\arraybackslash}X>{\centering\arraybackslash}X>{\centering\arraybackslash}X>{\centering\arraybackslash}X}
\phasecard{1}{milcVerde}{milcGris}{Escucha\\[-1mm]\small Observo, escucho y pregunto.} &
\phasecard{2}{milcTurquesa}{milcGris}{Entender\\[-1mm]\small Sistematizo y ordeno.} &
\phasecard{3}{milcMagenta}{milcGris}{Hacer\\[-1mm]\small Construyo y aplico.} &
\phasecard{4}{milcVino}{milcGris}{Cierre\\[-1mm]\small Evalúo y reflexiono.}
\end{tabularx}


\puente{Empezamos con un juego mental: imagina que un primo lejano te busca en internet por primera vez.}

\milcestacion{milcVerde}{Estación 1 de 3 · Escucha}

\begin{softbox}{milcVerde}{milcGris}{Escena para escuchar}
\textbf{Actividad 1 · IDENTIFICA --- Imagina al primo curioso} (10 min · individual). Cierra los ojos un momento. Imagina esto: un primo tuyo de otra ciudad, al que casi no conoces, escucha que estás en Cartago y \textbf{te busca en internet}. Ahora abre los ojos y \textbf{escribe en el cuaderno} qué cosas crees que encontraría sobre ti: \emph{¿Una foto? ¿Tu nombre? ¿Algún video tuyo? ¿Comentarios en juegos? ¿Tu Tik Tok? ¿Nada?}. \textbf{Anota 5 cosas} que crees que tu primo encontraría. No tienen que ser todas verdaderas; pueden ser cosas que tú \emph{sospechas} que hay.
\end{softbox}

{\sffamily\bfseries\fontsize{8}{10}\selectfont\textcolor{milcNegro!55}{MARCA CUANDO LO HAYAS HECHO}}
\milccheckrow{Cerré los ojos y pensé en el primo.}
\milccheckrow{Escribí 5 cosas que él encontraría de mí.}
\milccheckrow{No me dio pena escribir las verdaderas.}

\begin{infoband}{milcVerde}


\textbf{Lo que va al cuaderno (Actividad 1 · IDENTIFICA).} Título: «Actividad 1 --- El primo curioso». \textbf{Formato:} lista numerada de 5 cosas. \textbf{Extensión:} 5 viñetas cortas, una línea cada una. \textbf{Sabes que terminaste cuando} tienes 5 cosas escritas y ninguna te dio risa o pena al escribirla.
\end{infoband}

\puente{Después de pensar en el primo, vas a aprender 4 reglas para cuidar tu identidad digital sin esfuerzo.}

\milcestacion{milcTurquesa}{Estación 2 de 3 · Entender}

\begin{softbox}{milcTurquesa}{milcGris}{Lo que vas a entender}
Lo que acabaste de hacer es mirar tu identidad digital actual. Ahora vas a aprender \textbf{4 reglas} para cuidarla. Son simples y van a servirte toda la vida.
\end{softbox}

\milcpaso{1}{milcTurquesa}{\textbf{Regla 1 --- Tu nombre completo no se regala.} Tu nombre y apellidos completos son tuyos. No los pongas en sitios desconocidos. Para juegos y apps usa un \emph{apodo}: \texttt{LaraDelValle12} es mejor que \texttt{Lara Cárdenas Martínez}.}
\milcpaso{2}{milcTurquesa}{\textbf{Regla 2 --- Las fotos quedan.} Una foto que subes hoy puede aparecer en 10 años. Antes de subir una foto: ¿la mostrarías en clase? Si la respuesta es no, no la subas.}
\milcpaso{3}{milcTurquesa}{\textbf{Regla 3 --- Lo que escribes te firma.} Cada comentario, cada mensaje, lleva tu nombre o apodo. Si vas a escribir algo, léelo dos veces y pregunta: \emph{``¿Esto suena a mí o suena raro?''}.}
\milcpaso{4}{milcTurquesa}{\textbf{Regla 4 --- Lo importante NO se cuenta en internet.} Tu dirección, tu colegio exacto, los horarios que llegas a casa, fotos donde se ve el uniforme: esas cosas se cuentan en persona, a quien conoces de verdad.}

\begin{drawbox}{milcTurquesa}{Las 4 reglas para recordar}
\textbf{Regla 1.} Mi nombre completo no se regala. \\[1mm]
\textbf{Regla 2.} Las fotos quedan. \\[1mm]
\textbf{Regla 3.} Lo que escribo me firma. \\[1mm]
\textbf{Regla 4.} Lo importante se cuenta en persona.
\end{drawbox}

\begin{softbox}{milcMostaza}{milcGris}{Errores comunes y su corrección}
\textbf{Errores que cometen tus compañeros (y a veces tú también).} (1) Usar el nombre completo en juegos en línea --- los adultos malos buscan así. (2) Subir fotos donde se ve el uniforme del colegio --- cualquiera sabe dónde estás cada día. (3) Compartir contraseña con \emph{``mi mejor amigo''} --- la amistad cambia, la contraseña queda. (4) Aceptar solicitudes de gente que no conoces porque ``tienen muchos amigos en común''.
\end{softbox}

\begin{infoband}{milcTurquesa}


\textbf{Lo que va al cuaderno (Actividad 2 · EXPLICA).} Título: «Actividad 2 --- Las 4 reglas con mis palabras». \textbf{Formato:} las 4 reglas re-escritas \emph{con tu propia voz}, una por línea. Después dibuja 1 ícono al lado de cada una (lo que quieras). \textbf{Extensión:} 4 frases cortas + 4 íconos. \textbf{Sabes que terminaste cuando} podrías explicarle las 4 reglas a un hermano menor sin abrir el cuaderno.
\end{infoband}

\puente{Con las 4 reglas, dibujas tu propia tarjeta de identidad en el cuaderno.}

\milcestacion{milcMagenta}{Estación 3 de 3 · Hacer}

\begin{softbox}{milcMagenta}{milcGris}{Cómo vas a construir tu producto}
\textbf{Actividad 3 · CREA --- Mi tarjeta de identidad digital} (30 min · individual). Hora de hacer la tarjeta. Vas a usar \textbf{una hoja del cuaderno completa}, no media. Vas a dibujar, escribir y firmar. Sigue los 6 pasos con calma. No hay prisa.
\end{softbox}

\milcpaso{1}{milcMagenta}{\textbf{Paso 1.} En la parte superior de la hoja escribe en grande: \textbf{``Mi identidad digital''}. Debajo, tu nombre real (porque esta tarjeta es para ti, no para subir a internet).}
\milcpaso{2}{milcMagenta}{\textbf{Paso 2.} Dibuja un cuadro grande en el centro de la hoja. Adentro dibuja un \textbf{retrato tuyo} (un dibujo, no foto). No tiene que estar perfecto. La idea es que se parezca a ti, no a alguien más.}
\milcpaso{3}{milcMagenta}{\textbf{Paso 3.} Al lado del retrato, escribe \textbf{3 cosas que te gusta hacer}. Pueden ser cualquier cosa: jugar fútbol, dibujar mangas, oír reguetón, cocinar con tu abuela, hablar con tu hermana. Cosas reales, no las que ``deberías'' decir.}
\milcpaso{4}{milcMagenta}{\textbf{Paso 4.} Debajo, escribe \textbf{2 cosas que cuidas}. Son cosas que NO le cuentas a desconocidos. Ejemplo: tu dirección, una historia familiar, una contraseña, un secreto pequeño. No las escribas explícitas en la tarjeta, escribe solo \textbf{el tema}: \emph{``mi dirección''}, \emph{``mi familia''}.}
\milcpaso{5}{milcMagenta}{\textbf{Paso 5.} Abajo, escribe \textbf{1 frase} que te describa. Una frase que tú mismo inventes, no copiada de TikTok. Ejemplo: \emph{``Soy Lara, me gustan los gatos y las matemáticas raras.''}.}
\milcpaso{6}{milcMagenta}{\textbf{Paso 6.} En la esquina inferior derecha, firma con \textbf{tu nombre completo y la fecha de hoy}. Esta tarjeta queda firmada como tuya. La vas a pegar en la primera página del cuaderno de Tecnología.}

\begin{drawbox}{milcMagenta}{Antes de pegar la tarjeta}
\checkbox La tarjeta cabe en una hoja entera del cuaderno. \\[1mm]
\checkbox Hay un retrato tuyo dibujado. \\[1mm]
\checkbox Están 3 cosas que te gusta hacer (cosas reales). \\[1mm]
\checkbox Están 2 cosas que cuidas (solo el tema, no el contenido). \\[1mm]
\checkbox La frase que te describe la escribiste tú. \\[1mm]
\checkbox Firmaste con nombre y fecha.
\end{drawbox}

\begin{softbox}{milcMagenta}{milcGris}{Plantilla de trabajo}
\textbf{Estructura de la tarjeta.} \textit{Arriba:} título + tu nombre. \textit{Centro:} retrato dibujado. \textit{Lado derecho:} 3 cosas que te gustan. \textit{Abajo:} 2 cosas que cuidas + 1 frase tuya. \textit{Esquina:} firma + fecha.
\end{softbox}

\begin{infoband}{milcMagenta}


\textbf{Lo que va al cuaderno (Actividad 3 · CREA).} Título: «Actividad 3 --- Mi tarjeta de identidad digital». \textbf{Formato:} dibujo de hoja completa con texto y retrato. \textbf{Extensión:} 1 página entera del cuaderno. \textbf{Sabes que terminaste cuando} muestras la tarjeta al profe y él dice ``esa eres tú''.
\end{infoband}

\puente{El producto es esa tarjeta: la pegas en la primera página del cuaderno de Tecnología.}

\milcestacion{milcMostaza}{Producto de la sesión}
\textbf{Tarjeta de identidad digital · firmada y pegada en el cuaderno}.
\begin{softbox}{milcMostaza}{milcGris}{Criterios de calidad}
(1) La tarjeta ocupa una hoja completa. (2) Tiene los 5 elementos (título, retrato, gustos, cuidados, frase). (3) La frase que describe al estudiante fue inventada por él/ella. (4) Está firmada con nombre y fecha. (5) Cumple las 4 reglas aprendidas hoy (no datos sensibles visibles).
\end{softbox}

\puente{Antes de terminar, miras tu tarjeta y te preguntas: \emph{``Si mi mamá viera esto, ¿le gustaría cómo me presento?''}}

\milcestacion{milcVino}{Evaluación liberadora · cinco dimensiones}

\begin{softbox}{milcVino}{milcGris}{Sentido de la evaluación}
Evaluar no es solo revisar una nota. Es reconocer qué aprendí, qué cambió en mí, qué emociones manejé, cómo me relaciono con la ciudad y la comunidad, y qué le dejaré a quien viene detrás.
\end{softbox}

\textbf{Mi reflexión liberadora en cinco dimensiones.} Completa cada frase iniciada:

\small
\begin{tabularx}{\textwidth}{>{\bfseries\RaggedRight\arraybackslash}p{3.4cm}Y>{\RaggedRight\arraybackslash}p{4.6cm}}
\rowcolor{milcVino}\color{white}Dimensión & \color{white}\bfseries Mi respuesta (frame starter) & \color{white}\bfseries Algo que descubrí\\
1. Personal & Lo que más cambió en mí fue \linead{6.5cm} & \linead{4.2cm}\\
2. Emocional & La emoción más fuerte fue \linead{2.5cm} y la regulé \linead{3cm} & \linead{4.2cm}\\
3. Ciudadana & En democracia esto importa porque \linead{6cm} & \linead{4.2cm}\\
4. Local (Cartago) & En mi barrio sería útil para \linead{6.5cm} & \linead{4.2cm}\\
5. Intergeneracional & A mi abuela/abuelo le diría \linead{6.5cm} & \linead{4.2cm}\\
\end{tabularx}
\normalsize

\begin{infoband}{milcVino}
\textbf{Para la próxima sustentación.} Vuelve a esta tabla en seis meses. Verás que las respuestas cambian — no porque lo aprendido cambie, sino porque tú cambias. La evaluación liberadora es una conversación contigo a través del tiempo.
\end{infoband}


\puente{Y para cerrar, 3 voces antiguas tienen algo que decirte sobre cuidar la palabra propia.}

\milcestacion{milcGradeDark}{Triángulo de pensamiento · cierre filosófico}

\begin{softbox}{milcVino}{milcGris}{Dussel · lente del nosotros}
\textit{Mi voz cuenta tanto como cualquier otra. Mi nombre vale lo mismo que cualquier nombre.}

{\footnotesize\itshape\color{milcNegro!70}--- formulación de esta guía, al modo de Enrique Dussel. No es una cita textual.}

Hay un montón de gente en internet con nombres famosos, con miles de seguidores, con voces muy ruidosas. Pero \textbf{tu voz también cuenta}. Tu identidad digital no necesita parecerse a la de un youtuber para tener valor. Vale por ser tuya. Cuídala como vale.

\textbf{¿Estoy intentando parecer alguien más en internet, o estoy contando quién soy de verdad?}
\end{softbox}
\linead{\linewidth}\\[1mm]
\linead{\linewidth}

\begin{softbox}{milcMostaza}{milcGris}{Estoicismo · Marco Aurelio · cuidado interior}
\textit{Lo que escribes hoy queda. Escribe entonces como si lo fueras a leer dentro de 10 años.}

{\footnotesize\itshape\color{milcNegro!70}--- formulación de esta guía, al modo de Marco Aurelio. No es una cita textual.}

Hay cosas que escribes hoy de las que dentro de 10 años te vas a arrepentir. Otras te van a gustar leerlas. La diferencia está en cuidar lo que escribes, no en cuántos likes te dan ahora. Antes de subir algo, pregúntate: \emph{``¿Yo en 10 años, leyendo esto, me sentiría orgulloso?''}.

\textbf{¿Lo que voy a escribir aguanta el paso del tiempo, o me daría pena dentro de unos años?}
\end{softbox}
\linead{\linewidth}\\[1mm]
\linead{\linewidth}

\begin{softbox}{milcTurquesa}{milcGris}{Floridi · lente de la infoesfera}
\textit{Internet es una conversación enorme entre muchas voces. Tú también eres una de esas voces.}

{\footnotesize\itshape\color{milcNegro!70}--- formulación de esta guía, al modo de Luciano Floridi. No es una cita textual.}

Cada vez que escribes algo en internet, lo metes a una conversación gigante donde millones de personas hablan. Tu identidad digital es \textbf{tu manera de meterte a esa conversación}. Puedes meterte con respeto o gritando. Las dos formas cuentan, pero hacen mundos distintos.

\textbf{Cuando me meto a internet, ¿estoy aportando algo o solo agregando ruido?}
\end{softbox}
\linead{\linewidth}\\[1mm]
\linead{\linewidth}

\begin{infoband}{milcNegro}
\textbf{Nota para el docente.} Las tres formulaciones son del autor de la guía, no citas textuales de los pensadores: recogen su lente para pensar el tema, no sus palabras. Si vas a citarlos en clase, remite a la obra original.
\end{infoband}

\milcestacion{milcVino}{Mi autoevaluación · marca una casilla por fila}

{\small
\setlength{\extrarowheight}{2.2pt}
\begin{tabularx}{\textwidth}{>{\RaggedRight\arraybackslash\bfseries}p{3.0cm}YYY}
\rowcolor{milcVino}\color{white}\bfseries Criterio & \color{white}\bfseries Lo logré & \color{white}\bfseries En proceso & \color{white}\bfseries Necesito apoyo\\[.6mm]
Comprensión del tema & \milcopcion{Explico las ideas con mis palabras.} & \milcopcion{Reconozco las ideas, falta precisión.} & \milcopcion{Necesito repasar lo básico.}\\[.8mm]
\rowcolor{milcGris}Pensamiento computacional & \milcopcion{Usé los pilares al armar mi producto.} & \milcopcion{Usé algunos pilares.} & \milcopcion{Necesito releer la estación 2.}\\[.8mm]
Aplicación y producto & \milcopcion{Mi producto cumple los criterios.} & \milcopcion{Está armado, con ajustes pendientes.} & \milcopcion{Aún está en construcción.}\\[.8mm]
\rowcolor{milcGris}Reflexión 5 dimensiones & \milcopcion{Respondí cada una con honestidad.} & \milcopcion{Respondí algunas con detalle.} & \milcopcion{Mis respuestas son muy generales.}\\[.8mm]
Triángulo de pensamiento & \milcopcion{Conecté las 3 voces con mi trabajo.} & \milcopcion{Conecté al menos una.} & \milcopcion{No las relaciono con mi proceso.}\\[.8mm]
\end{tabularx}}

\vspace{3mm}
\textbf{Hoy entendí que:}\\[1mm]
\linead{\linewidth}\\[1mm]
\linead{\linewidth}

\textbf{Mi compromiso para la próxima sesión es:}\\[1mm]
\linead{\linewidth}\\[1mm]
\linead{\linewidth}

\begin{softbox}{milcMostaza}{milcGris}{Cita de despedida · Marco Aurelio}
\textit{``El obstáculo en el camino se vuelve el camino. Lo que estorba al camino se vuelve camino.''}\\[1mm]
Cada error de hoy, bien observado, se convierte en pieza del camino que pisarás mañana.
\end{softbox}

\small
\begin{tabularx}{\textwidth}{>{\bfseries}p{3.6cm}Y}
\rowcolor{milcGradeDark!12}Nombre completo: & \linead{10cm}\\
Fecha: & \linead{4cm} \quad Curso/grupo: \linead{4cm}\\
Mi firma: & \linead{9cm}\\
\end{tabularx}
\normalsize

\begin{infoband}{milcGradeDark}
\textbf{Para la próxima sesión.} Llega con tu \textbf{tarjeta pegada en la primera página del cuaderno}. La próxima clase activas tu correo institucional --- la primera dirección que firmas como tuya. Lo que aprendiste hoy te va a servir: el correo lleva tu nombre. Cuídalo como cuidas tu tarjeta.
\end{infoband}

\end{document}
